\section{Розділ~\ref{sec:dofaalt}. Інші теорії \nt{DOPA}}

Кожна з розширених картин спирається на свої прямі вимірювання дофаміну (розкид точок перевертання, відповіді на неприємне і нове, повільне зростання, відповіді на зміну смаку), але формули, які їх пояснюють, --- теорії, і щодо двох із них досліди поки суперечать одне одному.

{\footnotesize
\setlength{\tabcolsep}{3pt}\renewcommand{\arraystretch}{1.15}
\begin{longtable}{>{\raggedright}p{0.5cm}>{\raggedright}p{4.0cm}>{\raggedright}p{5.6cm}>{\raggedright}p{2.3cm}>{\raggedright\arraybackslash}p{1.7cm}}
\toprule
№ & Твердження & Дослід і що виміряно & Джерело & Статус \\
\midrule\endfirsthead
\toprule
№ & Твердження & Дослід і що виміряно & Джерело & Статус \\
\midrule\endhead
1 & Асиметричне правило~\eqref{eq:asymmetric} сходиться до точки~\eqref{eq:expectile}; при $\kappa=1/2$ це середнє, для краплі з імовірністю $1/2$ --- $V=\kappa$ (рис.~\ref{fig:expectiles}) & Точка~\eqref{eq:expectile} --- експектиль, розв'язок задачі найменших квадратів із різними вагами додатних і від'ємних відхилень; для прикладу розділу виведено в самому розділі. & \cite{NeweyPowell1987}; виведено в розділі & теорія \\
2 & У різних \nt{DOPA}-нейронів різні точки перевертання, і вони впорядковані за асиметрією (твердження~\ref{prop:expectile}) & Миші, оптогенетично мічені \nt{DOPA}-нейрони вентральної покришки, винагороди різного розміру: точка, де сплеск змінюється провалом, у різних нейронів різна; у нейронів із більшою асиметрією відповідей вона вища; за відповідями групи відновлюється форма розподілу винагород. Що це головна функція розкиду --- гіпотеза. & \cite{Dabney2020} & виміряно \\
3 & Частина \nt{DOPA}-нейронів відповідає на неприємне сплеском & Мавпи: одні \nt{DOPA}-нейрони збуджуються на провісник винагороди і гальмуються на провісник струменя повітря в око, інші збуджуються на обидва; групи розташовані в різних частинах середнього мозку. & \cite{Matsumoto2009, BrombergMartin2010} & виміряно \\
4 & Модуль помилки або новизна задає швидкість навчання & Прямого досліду, де сигнал важливості в \nt{DOPA}-нейронів змінює швидкість навчання, у цитованих роботах немає; огляд пропонує роль сигналу «увага, важливо». & \cite{BrombergMartin2010} & гіпотеза \\
5 & Окремий провід для осі небезпеки & Миші: \nt{DOPA}-нейрони, що йдуть до заднього кінця стріатума, відповідають на нові й сильні подразники, а не на винагороду; їхнє руйнування послаблює уникання нового предмета. & \cite{Menegas2018} & виміряно \\
6 & Оптимальний час дії $\tau_a^*=\sqrt{C/\bar R}$~\eqref{eq:vigor} & Мінімум суми $C/\tau_a+\bar R\tau_a$; виведено в розділі. У вихідній моделі швидкість дій задає ціна часу, що дорівнює середній винагороді. & \cite{Niv2007}; виведено в розділі & теорія \\
7 & Повільний рівень дофаміну несе $\bar R$ і задає швидкість дій (твердження~\ref{prop:multiplex}) & Щури, мікродіаліз і вольтамперометрія в прилеглому ядрі: рівень дофаміну від хвилини до хвилини йде за частотою винагород і за швидкістю дій. У людей L-DOPA посилює залежність часу реакції від середньої частоти винагород. Що повільний рівень дорівнює саме $\bar R$, не доведено. & \cite{Hamid2016, Beierholm2013vigor} & гіпотеза \\
8 & Повільний рівень складається з двох джерел: викид змінюється біля виходів без зміни частоти імпульсів, але повільне зростання буває і в самій частоті & Щури, одна задача: викид у прилеглому ядрі йде за змінним очікуванням винагороди, а частота імпульсів упізнаних \nt{DOPA}-нейронів вентральної покришки --- ні. Миші у віртуальному коридорі (рядок~10): плавне зростання під час наближення до винагороди є і в імпульсах упізнаних \nt{DOPA}-нейронів. & \cite{Mohebi2019, Kim2020} & виміряно \\
9 & Дофамін плавно росте, поки тварина наближається до далекої винагороди & Щури в лабіринті, вольтамперометрія в стріатумі: концентрація дофаміну наростає за секунди в міру наближення до мети, крутизна залежить від відстані і розміру винагороди. & \cite{Howe2013ramp, Hamid2016} & виміряно \\
10 & Зростання --- помилка $\delta$ при уточненні оцінки~\eqref{eq:ramp}; стрибок у коридорі дає сплеск (рис.~\ref{fig:ramp}) & Формула і число $\delta=0.75\,V$ на секунду --- виведено в розділі. Дослід: миші у віртуальному коридорі, миттєве перенесення ближче до мети; імпульси тіл, кальцій тіл і виходів та концентрація в стріатумі відповідають як помилка, а не як очікування $V$. Яке з двох пояснень правильне на всіх виходах, обговорюється. & \cite{Kim2020} & виміряно \\
11 & Погляд назад: дофамін повідомляє міру причинного зв'язку~\eqref{eq:retro} & Миші, викид дофаміну в прилеглому ядрі в дослідах, де ця теорія і помилка~\eqref{eq:ctd} передбачають різне: у низці дослідів викид ішов за першою. & \cite{Jeong2022} & гіпотеза \\
12 & Відповідь дофаміну під час навчання поступово зсувається від винагороди до провісника & Миші, сигнал виходів \nt{DOPA}-нейронів у черевному стріатумі в ході навчання: сплеск поступово переходить з моменту винагороди на момент провісника, як вимагає навчання за~\eqref{eq:ctd}. & \cite{Amo2022} & виміряно \\
13 & \nt{DOPA}-нейрони відповідають на зміну смаку винагороди за тієї самої цінності & Щури, запис нейронів вентральної покришки: заміна соку іншим, рівним за цінністю, спричиняє відповідь, хоча помилка~\eqref{eq:ctd} дорівнює нулю. & \cite{Takahashi2017} & виміряно \\
14 & Штучні сплески вчать зв'язку двох нейтральних подразників & Щури, дослід із попереднім поєднанням двох подразників без винагороди: оптогенетичне ввімкнення \nt{DOPA}-нейронів на переході від першого до другого створює зв'язок, вимкнення заважає його вивчити. & \cite{Sharpe2017} & виміряно \\
15 & Вектор помилок за ознаками~\eqref{eq:vectortd} & Теорія помилки передбачення ознак; прямої перевірки векторної форми немає. & \cite{Gardner2018} & гіпотеза \\
16 & Різні \nt{DOPA}-нейрони в одній задачі несуть різні величини & Миші у віртуальному коридорі, двофотонна зйомка \nt{DOPA}-нейронів: одні пов'язані з винагородою, інші зі швидкістю і поворотом, треті з провісниками і точністю вибору. & \cite{Engelhard2019} & виміряно \\
17 & Джгут замість проводу (твердження~\ref{prop:bundle}) & Зведення рядків~2--16: кожне розкладання виміряно окремо; єдиної картини, що всі вони --- частини одного сигналу, дослідом не встановлено. & --- & гіпотеза \\
\bottomrule
\end{longtable}}
